%% LyX 2.6.0~devel created this file.  For more info, see https://www.lyx.org/.
%% Do not edit unless you really know what you are doing.
\documentclass{article}
\usepackage{amsmath}
\usepackage{amssymb}
\usepackage{fontspec}
\setlength{\parindent}{0bp}
\usepackage{xcolor}
\definecolor{lyxboxbgcolor}{rgb}{0.800781, 0.800781, 0.800781}
\usepackage{float}
\usepackage{cancel}
\usepackage{geometry}
\geometry{verbose,tmargin=2.7cm,bmargin=2cm,lmargin=2cm,rmargin=2cm,headheight=2cm,headsep=0cm,footskip=0cm}
\usepackage{setspace}
\usepackage{esint}
\usepackage{microtype}
\onehalfspacing
\usepackage[pdfusetitle,
 bookmarks=true,bookmarksnumbered=true,bookmarksopen=true,bookmarksopenlevel=2,
 breaklinks=true,pdfborder={0 0 0},pdfborderstyle={},backref=section,colorlinks=true]
 {hyperref}
\hypersetup{
 linkcolor=black, citecolor=black}

\makeatletter

%%%%%%%%%%%%%%%%%%%%%%%%%%%%%% LyX specific LaTeX commands.
\providecommand\textquotedblplain{%
  \bgroup\addfontfeatures{Mapping=}\char34\egroup}

\@ifundefined{date}{}{\date{}}
%%%%%%%%%%%%%%%%%%%%%%%%%%%%%% User specified LaTeX commands.
% Engine guard: use XeLaTeX/LuaLaTeX
\usepackage{iftex}
\ifPDFTeX
  \PackageError{template}{Please compile with XeLaTeX or LuaLaTeX}{}
\fi

% ======== ADVANCED HEBREW ACADEMIC TEMPLATE ========
% Packages (trimmed to avoid conflicts)
% \usepackage{autobreak}  % REMOVE: can conflict with amsmath; use \allowdisplaybreaks
% \usepackage[T1]{fontenc}  % REMOVE: not needed with fontspec
\usepackage{relsize}
\usepackage{xcolor}
\usepackage{amssymb}
\usepackage{titlesec}
\usepackage{polyglossia}
\usepackage{fancyhdr}
\usepackage{tcolorbox}
\usepackage{tikz}
\usepackage{pgfplots}
\usepackage{booktabs}
\usepackage{array}
\usepackage{multirow}
\usepackage{enumitem}
\usepackage{float}
\usepackage{caption}
\usepackage{subcaption}
\usepackage{geometry}
\usepackage{microtype}
\usepackage{fontspec}
\usepackage{physics}
\usepackage{siunitx}
\usepackage{mathtools}
\usepackage{empheq}
\usepackage{mdframed}

% Colors (deduplicated)
\definecolor{primaryblue}{RGB}{33,150,243}
\definecolor{secondaryblue}{RGB}{63,81,181}
\definecolor{accentgreen}{RGB}{76,175,80}
\definecolor{warningcolor}{RGB}{255,152,0}
\definecolor{errorcolor}{RGB}{244,67,54}
\definecolor{infocolor}{RGB}{3,169,244}
\definecolor{notecolor}{RGB}{255,193,7}
\definecolor{theoremcolor}{RGB}{156,39,176}
\definecolor{definitioncolor}{RGB}{255,87,34}
\definecolor{examplecolor}{RGB}{0,150,136}
\definecolor{lightgray}{RGB}{245,245,245}
\definecolor{darkgray}{RGB}{66,66,66}
% Add aliases for colors referenced in layout to avoid missing-color errors
\colorlet{successcolor}{accentgreen}
\colorlet{warningred}{errorcolor}
\colorlet{purple}{theoremcolor}
\colorlet{green}{accentgreen}
\colorlet{blue}{primaryblue}
\colorlet{orange}{warningcolor}
\definecolor{brown}{RGB}{121,85,72}

% Language + fonts (polyglossia)
\setdefaultlanguage{hebrew}
\setotherlanguage{english}

% Robust font selection with fallbacks (fixes Times New Roman not loadable)
% Hebrew serif
\IfFontExistsTF{David CLM}
  {\newfontfamily\hebrewfont[Script=Hebrew]{David CLM}}
  {\IfFontExistsTF{Noto Serif Hebrew}
     {\newfontfamily\hebrewfont[Script=Hebrew]{Noto Serif Hebrew}}
     {\newfontfamily\hebrewfont[Script=Hebrew]{Noto Sans Hebrew}}}

% Hebrew sans
\IfFontExistsTF{Nachlieli CLM}
  {\newfontfamily\hebrewfontsf[Script=Hebrew]{Nachlieli CLM}}
  {\IfFontExistsTF{Noto Sans Hebrew}
     {\newfontfamily\hebrewfontsf[Script=Hebrew]{Noto Sans Hebrew}}
     {\newfontfamily\hebrewfontsf[Script=Hebrew]{David CLM}}}

% Latin main font (fallbacks for missing Times New Roman)
\IfFontExistsTF{Times New Roman}
  {\setmainfont{Times New Roman}}
  {\IfFontExistsTF{TeX Gyre Termes}
     {\setmainfont{TeX Gyre Termes}}
     {\IfFontExistsTF{Liberation Serif}
        {\setmainfont{Liberation Serif}}
        {\IfFontExistsTF{FreeSerif}
           {\setmainfont{FreeSerif}}
           {\setmainfont{Latin Modern Roman}}}}}

% siunitx
\sisetup{
  output-decimal-marker = {.},
  group-separator = {,},
  group-minimum-digits = 4,
  per-mode = symbol
}

% Geometry (avoid overriding in LyX header)
\geometry{
  a4paper,
  left=2cm, right=2cm, top=2.7cm, bottom=2cm,
  headheight=16pt, headsep=0.5cm, footskip=1cm
}

% Display math spacing (fix duplicates)
\AtBeginDocument{\setlength\abovedisplayskip{6pt}}
\AtBeginDocument{\setlength\abovedisplayshortskip{6pt}}
\AtBeginDocument{\setlength\belowdisplayskip{6pt}}
\AtBeginDocument{\setlength\belowdisplayshortskip{6pt}}

% Headers
\pagestyle{fancy}
\fancyhead{}
\fancyhead[R]{\nouppercase{\leftmark}}
\renewcommand{\headrulewidth}{0.5pt}

% Footnote rule (single definition)
\AtBeginDocument{
  \renewcommand\footnoterule{%
    \kern 3pt
    \hbox to \textwidth{\hfill\vrule height 0.5pt width 0.4\textwidth}
    \kern 4pt
  }%
}

% Word spacing
\spaceskip=1.3\fontdimen2\font plus 1\fontdimen3\font minus 1.5\fontdimen4\font

% Colors (deduplicated)
\definecolor{primaryblue}{RGB}{33,150,243}
\definecolor{secondaryblue}{RGB}{63,81,181}
\definecolor{accentgreen}{RGB}{76,175,80}
\definecolor{warningcolor}{RGB}{255,152,0}
\definecolor{errorcolor}{RGB}{244,67,54}
\definecolor{infocolor}{RGB}{3,169,244}
\definecolor{notecolor}{RGB}{255,193,7}
\definecolor{theoremcolor}{RGB}{156,39,176}
\definecolor{definitioncolor}{RGB}{255,87,34}
\definecolor{examplecolor}{RGB}{0,150,136}
\definecolor{lightgray}{RGB}{245,245,245}
\definecolor{darkgray}{RGB}{66,66,66}

% QED symbol
%\renewcommand{\qedsymbol}{$\blacksquare$}

% Section formatting (wrap \@ macro with makeatletter)
\makeatletter
\renewcommand*{\@seccntformat}[1]{\hspace{0.5cm}\csname the#1\endcsname\hspace{0.5cm}}
\makeatother
\titleformat{\section}{\fontsize{20}{20}\bfseries}{\thesection}{10pt}{}
\titleformat{\subsection}{\fontsize{15}{15}\bfseries}{\thesubsection}{10pt}{}
\titleformat{\subsubsection}{\bfseries}{\thesubsubsection}{10pt}{}

% RTL equation tag direction
\makeatletter
\def\maketag@@@#1{\hbox{\m@th\normalfont\LRE{#1}}}
\def\tagform@#1{\maketag@@@{(\ignorespaces#1\unskip)}}
\makeatother

% Disjoint union, large logical and
\makeatletter
\def\moverlay{\mathpalette\mov@rlay}
\def\mov@rlay#1#2{\leavevmode\vtop{%
   \baselineskip\z@skip \lineskiplimit-\maxdimen
   \ialign{\hfil$\m@th#1##$\hfil\cr#2\crcr}}}
\newcommand{\charfusion}[3][\mathord]{#1{\ifx#1\mathop\vphantom{#2}\fi\mathpalette\mov@rlay{#2\cr#3}}\ifx#1\mathop\expandafter\displaylimits\fi}
\makeatother
\newcommand{\cupdot}{\charfusion[\mathbin]{\cup}{\cdot}}
\newcommand{\bigcupdot}{\charfusion[\mathop]{\bigcup}{\cdot}}
\newcommand{\kaliLargeLand}{\mathop{\mathlarger{\mathlarger{\mathlarger{\land}}}}}

% Hebrew text helper (fix)
\newcommand{\hebtext}[1]{\texthebrew{#1}}

% Define theorem-like environments
\newtheorem{theorem}{משפט}[section]
\newtheorem{lemma}[theorem]{למה}
\newtheorem{proposition}[theorem]{הצעה}
\newtheorem{corollary}[theorem]{מסקנה}
\newtheorem{definition}[theorem]{הגדרה}
\newtheorem{example}[theorem]{דוגמה}
\newtheorem{exercise}[theorem]{תרגיל}
\newtheorem{problem}[theorem]{בעיה}
\newtheorem{remark}[theorem]{הערה}
\newtheorem{notation}[theorem]{סימון}
\newtheorem{formula}[theorem]{נוסחה}
\newtheorem{result}[theorem]{תוצאה}
\newtheorem{experiment}[theorem]{ניסוי}
\newtheorem{physicslaw}[theorem]{חוק פיזיקלי}

% Define additional physics environments
\newenvironment{physicsnote}{\begin{notebox}{}}{\end{notebox}}
\newenvironment{physicsconstant}{\begin{mainbox}{קבוע פיזיקלי}}{\end{mainbox}}
\newenvironment{physicsunits}{\begin{mainbox}{יחידות}}{\end{mainbox}}
\newenvironment{physicsapplication}{\begin{examplebox}{יישום}}{\end{examplebox}}
\newenvironment{important}{\begin{importantbox}{}}{\end{importantbox}}
\newenvironment{summary}{\begin{summarybox}{}}{\end{summarybox}}

% Define proof environment
\newenvironment{solution}{\begin{proof}[פתרון]}{\end{proof}}

% ======== TCOLORBOX CONFIGURATIONS ========
\tcbuselibrary{breakable,skins,theorems,fitting,xparse}

% Main content box
\newtcolorbox{mainbox}[2][]{
    enhanced,
    breakable,
    colback=lightgray!20,
    colframe=primaryblue,
    fonttitle=\bfseries\large,
    title={#2},
    attach boxed title to top center={yshift=-2mm},
    boxed title style={size=small, colback=primaryblue, sharp corners},
    rounded corners,
    left=6mm,
    right=6mm,
    top=4mm,
    bottom=4mm,
    drop shadow,
    #1
}

% Theorem box
\newtcolorbox{theorembox}[2][]{
    enhanced,
    breakable,
    colback=theoremcolor!10,
    colframe=theoremcolor,
    fonttitle=\bfseries\large,
    title={משפט: #2},
    attach boxed title to top left={yshift=-2mm, xshift=5mm},
    boxed title style={size=small, colback=theoremcolor, sharp corners},
    rounded corners,
    left=6mm,
    right=6mm,
    top=4mm,
    bottom=4mm,
    drop shadow,
    #1
}

% Definition box
\newtcolorbox{definitionbox}[2][]{
    enhanced,
    breakable,
    colback=definitioncolor!10,
    colframe=definitioncolor,
    fonttitle=\bfseries\large,
    title={הגדרה: #2},
    attach boxed title to top left={yshift=-2mm, xshift=5mm},
    boxed title style={size=small, colback=definitioncolor, sharp corners},
    rounded corners,
    left=6mm,
    right=6mm,
    top=4mm,
    bottom=4mm,
    drop shadow,
    #1
}

% Example box
\newtcolorbox{examplebox}[2][]{
    enhanced,
    breakable,
    colback=examplecolor!10,
    colframe=examplecolor,
    fonttitle=\bfseries\large,
    title={דוגמה: #2},
    attach boxed title to top right={yshift=-2mm, xshift=-5mm},
    boxed title style={size=small, colback=examplecolor, sharp corners},
    rounded corners,
    left=6mm,
    right=6mm,
    top=4mm,
    bottom=4mm,
    drop shadow,
    #1
}

% Problem box
\newtcolorbox{problembox}[2][]{
    enhanced,
    breakable,
    colback=primaryblue!10,
    colframe=primaryblue,
    fonttitle=\bfseries\large,
    title={בעיה: #2},
    attach boxed title to top right={yshift=-2mm, xshift=-5mm},
    boxed title style={size=small, colback=primaryblue, sharp corners},
    rounded corners,
    left=6mm,
    right=6mm,
    top=4mm,
    bottom=4mm,
    drop shadow,
    #1
}

% Solution box
\newtcolorbox{solutionbox}[1][]{
    enhanced,
    breakable,
    colback=successcolor!10,
    colframe=successcolor,
    fonttitle=\bfseries\large,
    title=פתרון,
    attach boxed title to top right={yshift=-2mm, xshift=-5mm},
    boxed title style={size=small, colback=successcolor, sharp corners},
    rounded corners,
    left=6mm,
    right=6mm,
    top=4mm,
    bottom=4mm,
    drop shadow,
    #1
}

% Note box
\newtcolorbox{notebox}[2][]{
    enhanced,
    breakable,
    colback=notecolor!15,
    colframe=notecolor,
    fonttitle=\bfseries,
    title={הערה: #2},
    attach boxed title to top left={yshift=-2mm, xshift=5mm},
    boxed title style={size=small, colback=notecolor, sharp corners},
    rounded corners,
    left=6mm,
    right=6mm,
    top=4mm,
    bottom=4mm,
    drop shadow,
    #1
}

% Warning box
\newtcolorbox{warningbox}[2][]{
    enhanced,
    breakable,
    colback=warningcolor!15,
    colframe=warningcolor,
    fonttitle=\bfseries,
    title={אזהרה: #2},
    attach boxed title to top center={yshift=-2mm},
    boxed title style={size=small, colback=warningcolor, sharp corners},
    rounded corners,
    left=6mm,
    right=6mm,
    top=4mm,
    bottom=4mm,
    drop shadow,
    #1
}

% Important box
\newtcolorbox{importantbox}[2][]{
    enhanced,
    breakable,
    colback=errorcolor!10,
    colframe=errorcolor,
    fonttitle=\bfseries,
    title={חשוב: #2},
    attach boxed title to top center={yshift=-2mm},
    boxed title style={size=small, colback=errorcolor, sharp corners},
    rounded corners,
    left=6mm,
    right=6mm,
    top=4mm,
    bottom=4mm,
    drop shadow,
    #1
}

% Summary box
\newtcolorbox{summarybox}[2][]{
    enhanced,
    breakable,
    colback=infocolor!10,
    colframe=infocolor,
    fonttitle=\bfseries\large,
    title={סיכום: #2},
    attach boxed title to top center={yshift=-2mm},
    boxed title style={size=small, colback=infocolor, sharp corners},
    rounded corners,
    left=6mm,
    right=6mm,
    top=4mm,
    bottom=4mm,
    drop shadow,
    #1
}

% Formula box
\newtcolorbox{formulabox}[2][]{
    enhanced,
    breakable,
    colback=purple!10,
    colframe=purple,
    fonttitle=\bfseries\large,
    title={נוסחה: #2},
    attach boxed title to top center={yshift=-2mm},
    boxed title style={size=small, colback=purple, sharp corners},
    rounded corners,
    left=6mm,
    right=6mm,
    top=4mm,
    bottom=4mm,
    drop shadow,
    #1
}

% Result box
\newtcolorbox{resultbox}[2][]{
    enhanced,
    breakable,
    colback=accentgreen!15,
    colframe=accentgreen,
    fonttitle=\bfseries\large,
    title={תוצאה: #2},
    attach boxed title to top center={yshift=-2mm},
    boxed title style={size=small, colback=accentgreen, sharp corners},
    rounded corners,
    left=6mm,
    right=6mm,
    top=4mm,
    bottom=4mm,
    drop shadow,
    #1
}

% Enhanced mathematical environments
\newenvironment{calculation}
  {\begin{align}}
  {\end{align}}

\newenvironment{steps}
  {\begin{enumerate}[label=שלב \arabic*:, rightmargin=1cm]}
  {\end{enumerate}}

% Custom commands for common elements
\newcommand{\highlight}[1]{\textcolor{primaryblue}{\textbf{#1}}}
% CHANGED: avoid clash with the 'important' environment
\newcommand{\imptext}[1]{\textcolor{errorcolor}{\textbf{#1}}}
\newcommand{\note}[1]{\textcolor{notecolor}{\textit{#1}}}

% Enhanced equation highlighting
\newcommand{\boxedeq}[1]{\boxed{\displaystyle #1}}
\newcommand{\redeq}[1]{\textcolor{errorcolor}{\boxed{\displaystyle #1}}}
\newcommand{\blueeq}[1]{\textcolor{primaryblue}{\boxed{\displaystyle #1}}}

\makeatother

\begin{document}
\title{סיכום הרצאה - משוואות (77313)}
\maketitle
\begin{center}
כותב: ערן ריחני
\par\end{center}

\begin{center}
סמסטר א' תשפ\textquotedblright ה, האוניברסיטה העברית
\par\end{center}

\newpage{}

\tableofcontents{}
\begin{center}
\vfill{}
\par\end{center}

\begin{center}
הערות
\par\end{center}

\[
***
\]

\begin{center}
\par\end{center}

\[
***
\]

\begin{center}
\par\end{center}

\newpage{}

\section{פונקציות מוכללות }

נתחיל מהפונקציה $\phi_{1}\left(x\right)$
\begin{enumerate}
\item יש לה מקסימום ב-$x=0$ 
\item שואפת ב-$\pm\infty$ ל-0 (זוגית)
\item $\intop^{\infty}_{-\infty}\phi_{1}\left(x\right)dx=1$
\end{enumerate}
לורנציאן :

\[
\phi^{\left(a\right)}_{1}\left(x\right)=\frac{1}{\pi}\frac{1}{1+x^{2}}
\]

גאוסיאן:
\[
\phi^{\left(b\right)}_{1}\left(x\right)=\frac{1}{\sqrt{\pi}}e^{-x^{2}}
\]

נראה שהאינטגרלים של הפונקציות האלו שווה 1 .

\[
\intop^{\infty}_{-\infty}\phi^{\left(a\right)}_{1}\left(x\right)dx=\frac{1}{\pi}\intop^{\infty}_{-\infty}\frac{1}{1+x^{2}}dx=\frac{1}{\pi}\arctan\left(x\right)\intop^{\infty}_{-\infty}\frac{1}{\pi}\left(\frac{\pi}{2}+\frac{\pi}{2}\right)=1
\]

כעת נבצע את האינטגרל על $\phi^{\left(b\right)}_{1}$:
\[
\intop^{\infty}_{-\infty}\phi^{\left(b\right)}_{1}\left(x\right)dx=\frac{1}{\sqrt{\pi}}\underset{I}{\underbrace{\intop^{\infty}_{-\infty}e^{-x^{2}}dx}}=
\]

צ\textquotedbl ל $I=\sqrt{\pi}$

\begin{align*}
I^{2} & =\intop^{\infty}_{-\infty}e^{-x^{2}}dx\intop^{\infty}_{-\infty}e^{-y^{2}}dy=\iint^{\infty}_{\infty}\exp-\left(x^{2}+y^{2}\right)dxdy\overset{x^{2}+y^{2}=r^{2}}{\underset{\frac{y}{x}=\tan\varphi}{=}}\intop^{2\pi}_{0}\intop^{\infty}_{0}e^{-r^{2}}drd\varphi\\
 & =2\pi\intop^{\infty}_{0}e^{-r^{2}}dr\overset{r^{2}=z}{=}2\frac{\pi}{2}\intop^{\infty}_{0}e^{z}dz=\pi e^{-z}\mid^{0}_{\infty}=\pi
\end{align*}

\begin{important}
אינטגרל של גאוסיאן 

\[
\intop^{\infty}_{-\infty}e^{-\beta x^{2}}dx=\sqrt{\frac{\pi}{\beta}}
\]
\end{important}
נגדיר סדרת פונקציות : $\phi_{m}\left(x\right)=m\phi_{1}\left(mx\right)$
כאשר $m=1,2,\ldots$ 

ככל ש-$m$ גדל כך מתכוונת ומוגברת . 
\[
\intop^{\infty}_{-\infty}\phi_{m}\left(x\right)dx=\intop^{\infty}_{-\infty}m\phi_{1}\left(x\right)dx\overset{y=mx}{\underset{dy=mdx}{=}}\intop^{\infty}_{-\infty}\phi_{1}\left(y\right)dy=1
\]

\[
\begin{cases}
\phi^{\left(a\right)}_{m}=\frac{m}{\pi}\frac{1}{1+\left(mx\right)^{2}}\\
\phi^{\left(b\right)}_{m}=\frac{m}{\sqrt{\pi}}\exp\left(-m^{2}x^{2}\right)
\end{cases}
\]

אלו הן פונקציות יוצרות של פונקציית $\delta$.

מה קורה כאשר $m\longrightarrow\infty$ ?
\[
\underset{x\neq0,\ m\rightarrow\infty}{\phi_{m}}\left(x\right)=m\underset{x\neq0,\ m\rightarrow\infty}{\phi_{1}\left(mx\right)}\rightarrow0
\]

\[
\lim\phi_{m}\left(x\right)=\lim m\phi_{1}\left(x\right)\leqslant\lim_{m\rightarrow\infty}m\cdot\frac{1}{\left(mx\right)^{\alpha}}=x^{-\alpha}m^{1-\alpha}\underset{m\rightarrow\infty}{\longrightarrow}0
\]

לפיכך נגדיר 
\[
\delta\left(x\right)=\lim_{m\rightarrow\infty}\phi_{m}\left(x\right)
\]
\[
\delta\left(x\right)=\begin{cases}
0 & x\neq0\\
\infty & x=0
\end{cases}
\]
\[
\intop^{\infty}_{-\infty}\delta\left(x\right)dx=1
\]

שימושים:
\begin{enumerate}
\item גדלים נקודתיים - 
\[
\varphi\left(x\right)=3\delta\left(x\right)\Rightarrow\intop^{\infty}_{-\infty}\varphi\left(x\right)dx=3\intop^{\infty}_{-\infty}\varphi\left(x\right)dx=3
\]

\begin{enumerate}
\item לדוגמה, צפיפות מטען/מסה 
\begin{align*}
\lambda\left(x\right) & =m\delta\left(x-a\right)\\
m & =\intop^{\infty}_{-\infty}\lambda\left(x\right)dx=m\intop^{\infty}_{-\infty}\delta\left(x-a\right)dx\underset{y=x-a}{=}m\intop^{\infty}_{-\infty}\delta\left(y\right)dy=m
\end{align*}
\end{enumerate}
\end{enumerate}
טענה: 
\[
\forall\varepsilon>0\ \intop^{\infty}_{-\infty}\delta\left(x\right)dx=\intop^{\varepsilon}_{-\varepsilon}\delta\left(x\right)dx=1
\]

זהות חשובה: 
\[
\intop^{\infty}_{-\infty}\delta\left(x-a\right)f\left(x\right)dx=f\left(a\right)
\]


\subsection*{זהויות חשובות של פונקציה דלתא }
\begin{enumerate}
\item $\delta\left(ax\right)=\frac{1}{}$
\item $\delta\left(f\left(x\right)\right)=$
\end{enumerate}
%
\newpage{}

\section{שבוע שני הרצאה 4 - פונקציות מוכללות 30 October 2025 }

תזכורת: 
\[
\delta\left(x\right)=\begin{cases}
0 & x\neq0\\
\infty & x=0
\end{cases},\quad\intop^{\infty}_{-\infty}\delta\left(x\right)=1,\ \ \intop^{\infty}_{-\infty}\delta\left(x-a\right)f\left(x\right)=f\left(a\right)
\]

\[
\Theta\left(x\right)=\intop^{\infty}_{-\infty}\delta\left(x'\right)dx'=\begin{cases}
0 & x<0\\
1 & x>0
\end{cases}
\]

ראינו שפונקציית $\Theta$ היא הגבול: 
\[
\Theta\left(x\right)=\lim_{m\rightarrow\infty}\frac{1}{\pi}\left[\arctan\left(mx\right)+\frac{\pi}{2}\right]
\]

והנגזרת של פונקציית דלתא $\Theta$ היא $\delta$ .
\[
\delta\left(x\right)=\Theta'\left(x\right)
\]

סוף התזכורת.

נחשב את האינטגרל הבא:
\begin{align*}
I=\intop^{\infty}_{-\infty}\delta\left(x\right)\Theta\left(x\right)dx & =\Theta\left(x\right)^{2}\Biggr|^{\infty}_{-\infty}-\intop^{\infty}_{-\infty}\delta\left(x\right)\Theta\left(x\right)dx=1-\underset{I}{\underbrace{\intop^{\infty}_{-\infty}\delta\left(x\right)\Theta\left(x\right)dx}}\\
\Rightarrow & 2I=1\Rightarrow I=\frac{1}{2}
\end{align*}

זה מוביל אותו אל המשפט הבא:

\[
\intop^{\infty}_{-\infty}\delta\left(x-a\right)f\left(x\right)dx=\frac{f\left(a\rightarrow a_{+}\right)-f\left(a-a_{-}\right)}{2}
\]
 כלומר אם $f$ לא רציפה בנקודה $a$ אז נמצע את ערכי הפונקציה בנקודת
אי הרציפות.

דוגמה:

\[
f\left(x\right)=\begin{cases}
x^{3} & 0<x<1\\
x^{2}+2 & x<1
\end{cases}
\]

נאיבית היינו רוצים לגזור כל תחום בנפרד אבל אז היינו נתקלים בבעיה שאנו
לא יודעים מה קורה בנקודת אי-הרציפות.

הפתרון הוא להגדיר פונקציה חדשה ולהשתמש בפונקציה מוכללת
\[
\varphi\left(x\right)=\begin{cases}
x^{3} & 0\leqslant x\leqslant1\\
x^{2} & x>1
\end{cases},\quad\varphi'\left(x\right)=\begin{cases}
3x^{2} & x<1\\
2x & x>1
\end{cases}
\]
\[
f\left(x\right)=\varphi\left(x\right)+2\Theta\left(x-1\right)
\]
 כעת ניתן לגזור את הפונקציה בשלמותה.
\[
f'\left(x\right)=\varphi'\left(x\right)+2\delta\left(x-1\right)
\]


\subsubsection*{נגזרת של פונקצצית $\delta$}

\[
\delta\left(x\right)=\lim_{m\rightarrow\infty}\phi^{\left(b\right)}_{m}\left(x\right)=\frac{m}{\sqrt{\pi}}e^{-\left(mx\right)^{2}}
\]
\begin{align*}
\delta'\left(x\right) & =\lim_{m\rightarrow\infty}\phi'^{\left(b\right)}_{m}\left(x\right)\\
 & =\frac{m}{\sqrt{\pi}}\frac{d}{dx}\big(e^{-(mx)^{2}}\big)\\
 & =\frac{d}{dx}e^{-(mx)^{2}}\\
 & =e^{-(mx)^{2}}\cdot(-2mx)
\end{align*}
\[
\phi_{m}'(x)=\frac{m}{\sqrt{\pi}}(-2mx)e^{-(mx)^{2}}=-\frac{2m^{2}x}{\sqrt{\pi}}e^{-(mx)^{2}}
\]
\begin{align*}
\delta''\left(x\right) & =\lim_{m\rightarrow\infty}\phi''^{\left(b\right)}_{m}\left(x\right)\\
 & =\frac{d}{dx}\left(-\frac{2m^{2}x}{\sqrt{\pi}}e^{-(mx)^{2}}\right)\\
 & =-\frac{2m^{2}}{\sqrt{\pi}}\left(e^{-(mx)^{2}}+x\frac{d}{dx}e^{-(mx)^{2}}\right)\\
 & =-\frac{2m^{2}}{\sqrt{\pi}}\left(e^{-(mx)^{2}}-2m^{2}x^{2}e^{-(mx)^{2}}\right)
\end{align*}
\[
\delta_{m}''(x)=\frac{2m^{2}}{\sqrt{\pi}}e^{-(mx)^{2}}(2m^{2}x^{2}-1)
\]

לאחר השוואה לאפס לקבלת נקודות הקיצון נקבל 
\[
\delta'\left(x^{^{\ast}}\right)=\sqrt{\frac{2}{\pi e}}m^{2},\quad\delta'\left(x^{^{\ast}}_{-}\right)=-\sqrt{\frac{2}{\pi e}}m^{2}
\]

\selectlanguage{english}%
\begin{center}
insert here tickz graph of $\delta'\left(x\right)$ 
\par\end{center}

\selectlanguage{hebrew}%
נשים לב שזה מזכיר לנו דיפול חשמלי.

כעת נסתכל על האינטגרל הבא:
\begin{align*}
\intop^{\infty}_{-\infty}f\left(x\right)\delta'\left(x-a\right)dx & =\cancel{f\left(x\right)\delta\left(x-a\right)\Biggr|^{\infty}_{-\infty}}-\intop f'\left(x\right)\delta\left(x-a\right)dx\\
 & =-f'\left(a\right)
\end{align*}

ניתן להכליל את זה לנגזרות גבוהות יותר.

הכנס כאן את הנוסחה הכללית.

מהו האינטגרל הבא:
\[
\intop^{\infty}_{-\infty}\delta'\left(x-a\right)dx=\delta\left(x\right)\Biggr|^{\infty}_{-\infty}=0
\]

לכן $\delta$ היא מועמדת מעולה לתאר צפיפות מטען של דיפול. (סה\textquotedbl כ
המטען הוא אפס).

\[
\rho\left(x\right)=P_{0}\delta'\left(x\right)
\]
\[
\intop^{\infty}_{-\infty}\rho_{dp}\left(x\right)dx=0,\quad\intop^{\infty}_{-\infty}x\rho_{dp}dx=-P_{0}
\]
קיבלנו שהמומנט שלילי כיוון שהמטען מכוון ממינוס לפלוס והוקטור מכוון
ככה.

ניתן להכליל את $\delta\left(x\right)$ למימדים יותר גבוהים.

\[
\delta\left(x\right)\delta\left(y\right)\delta\left(z\right)=\begin{cases}
\infty & x,y,z=0\\
0 & x,y,z\neq0
\end{cases}
\]
\[
\intop^{\infty}_{-\infty}\delta^{3}\left(z,y,z\right)=1,\quad\left[\delta^{3}\left(x,y,z\right)\right]=\frac{1}{m^{3}}
\]


\subsection*{טרנספורם פורייה של פונקציית $\delta$}

\[
\mathcal{F}\left(\omega\right)=\frac{1}{\sqrt{2\pi}}\intop^{\infty}_{-\infty}f\left(t\right)e^{i\omega t}dt
\]

והטרנספורם ההפוך
\[
f\left(t\right)=\frac{1}{\sqrt{2\pi}}\intop^{\infty}_{-\infty}F\left(\omega\right)e^{-i\omega t}d\omega
\]
\[
f\left(t\right)=\frac{1}{\sqrt{2\pi}}\intop^{\infty}_{-\infty}e^{i\omega t}\delta\left(t\right)d\omega=\frac{1}{\sqrt{2\pi}}
\]

טענה:
\[
\delta\left(t\right)=\frac{1}{2\pi}\intop^{\infty}_{-\infty}e^{i\omega t}d\omega
\]

הוכחה: 
\[
insert\ proof\ here
\]

נוכיח שמקבלים את הפונקציה המקורית ע\textquotedbl י האינטגרל הבא:
\[
f\left(t\right)=\frac{1}{\sqrt{2\pi}}\intop^{\infty}_{-\infty}F\left(\omega\right)e^{-i\omega t}d\omega
\]
\begin{align*}
f\left(t\right) & =\intop^{\infty}_{-\infty}\frac{1}{\sqrt{2\pi}}\underset{F\left(\omega\right)}{\underbrace{\intop^{\infty}_{-\infty}F\left(\omega\right)e^{-i\omega t'}dt'}}e^{-i\omega t}d\omega\\
 & =\frac{1}{2\pi}\intop^{\infty}_{-\infty}dt'f\left(t'\right)\underset{2\pi\delta\left(t-t'\right)}{\underbrace{\intop^{\infty}_{-\infty}d\omega e^{i\omega\left(t'-t\right)}}}=f\left(t\right)
\end{align*}


\section{חשבון ווריאציה }

היסטוריה: 

ברנולי התעניין בשאלה באיזה מסלול הזמן של נפילה של כדור יהיה מינימלי?

עד עכשיו דיברנו על פונקציות $f:\ \mathbb{R}\rightarrow\mathbb{R}$ 

כעת נדבר על פונקציונל שמקבל פונקציה ומוציאה מספר.

\[
T_{AB}=\intop^{B}_{A}\frac{ds}{v}
\]
\[
ds=\sqrt{dx^{2}+dy^{2}}=dx\sqrt{1+y'\left(x\right)^{2}}
\]

מצאנו את $ds$ כעת נמצא את $v$ מחוק שימור האנרגיה בין 2 נקודות.
\[
0=-mgy+\frac{1}{2}mv^{2}\Rightarrow v=\sqrt{2gy}
\]

\[
T_{AB}=\intop^{x_{0}}_{0}\frac{dx\sqrt{1+y'\left(x\right)^{2}}}{\sqrt{2gy}}=\frac{1}{\sqrt{2gy}}\intop^{x_{0}}_{0}\sqrt{\frac{1+y'^{2}}{y}}dx
\]

קיבלנו גדול שנקרא פונקציונל שנסמנו ב-$I\left[y\left(x\right)\right]$
\[
I\left[y\left(x\right)\right]=\intop^{b}_{a}F\left(x,y,y',\ldots\right)dx
\]
המטרה היא למצוא $y\left(x\right)$שממזער (או ממקסם) את $I\left[y\left(x\right)\right]$. 

כמעט כל חוק טבע ניתן לקבל מתוך חשבון וריאציות ועקרון מזעור אנרגיה.

עם זאת חוקים אמפירים כגון חוק אוהם לא ניתן לקבל מתוכם.
\end{document}
